\documentclass[11pt,a4paper]{article}
\usepackage[margin=2.3cm]{geometry}
\usepackage{amsmath,amssymb}
\usepackage{booktabs}
\usepackage{graphicx}
\usepackage{xcolor}
\usepackage[colorlinks=true,linkcolor=blue,urlcolor=blue,citecolor=blue]{hyperref}
\usepackage{siunitx}

\newcommand{\tauh}{\tau_{h}}
\newcommand{\tautau}{\tauh\tauh}
\newcommand{\mtautau}{m_{\tau\tau}}
\newcommand{\mvis}{m_{\mathrm{vis}}}
\newcommand{\mtau}{m_{\tau}}
\newcommand{\kl}{\kappa_{\lambda}}
\newcommand{\dR}{\Delta R}
\newcommand{\ptmiss}{p_{T}^{\mathrm{miss}}}
\newcommand{\yes}{\textcolor{teal}{$\checkmark$}}
\newcommand{\no}{\textcolor{red}{$\times$}}

\title{\vspace{-1.5cm}Delphes vs.\ CMS for HH$\rightarrow b\bar b\tau\tau$:\\
\large where we stand, and one decision to take}
\author{Jingjing Pan}
\date{\today}

\begin{document}
\maketitle

\begin{abstract}
\noindent
This note has three parts. \textbf{Section~1} audits what the Delphes fast simulation
currently does relative to CMS NanoAOD, object by object: what is matched, what is tuned
to a CMS anchor, what differences remain by choice or by limitation, and which of the ten
NSBI input features still disagree. \textbf{Section~2} states the one open decision ---
whether to continue correcting the jet-based $\tauh$ or to rebuild it from the Delphes
particle-flow candidates --- with the considerations on each side and the questions we
would like $\tau$ expertise on. \textbf{Section~3} is a reference list of the studies
behind the numbers, for anyone who wants to check a specific claim.

Short summary: the $\tautau$ yield deficit that motivated this work (a factor
$\sim$5) is closed, the cutflow now agrees with CMS at every stage ($1.01$), and eight of
the ten NSBI features agree within statistics. The two that do not are $\mtautau$
(residual $\sim$7\,\si{\GeV} offset) and $\dR_{\tau\tau}$ (an excess at $2.2$--$3.2$).
\end{abstract}

\section{Audit: what we do, relative to CMS}

The Delphes card is the stock CMS card plus the patches in
\texttt{docs/card\_patch\_validation.md}; all object tuning is applied \emph{downstream},
so the produced samples are never re-made. Every efficiency and scale is measured on a
private 2024 NanoAODv15 anchor and applied to Delphes, so ``tuned'' below always means
``measured on CMS and transferred''.

\subsection{Matched or tuned to CMS}

Table~\ref{tab:matched} lists the quantities that either agree with CMS already or are
transferred from it by a measured map. In every case the map is derived on the anchor and
applied to Delphes downstream; none of them requires a change to the card.

\begin{table}[htbp]
\centering\small
\begin{tabular}{llp{6.6cm}}
\toprule
quantity & status & how \\
\midrule
b-tag efficiency ($b$, $c$, light) & \yes\ tuned & stochastic re-tag from \texttt{Jet.Flavor} at the anchor-measured UParT-AK4 Medium efficiency \\
b-jet energy scale & \yes\ tuned & response inverted to unity against \texttt{GenJet} \\
$m_{bb}$ peak and width & \yes\ matched & follows from the above \\
$\tauh$ identification efficiency & \yes\ tuned & anchor \texttt{GenVisTau}$\rightarrow$DeepTau v2.5 Medium; includes CMS's HPS reconstruction inefficiency \\
$\tauh$ visible mass & \yes\ tuned & \emph{sampled} per jet from the anchor's per-$p_{T}$ mass quantiles \\
$\tauh$ energy scale & \yes\ tuned & CMS/Delphes response ratio, both against a \emph{visible} $\tau$ reference \\
$\tauh$ acceptance & \yes\ identical & common $p_{T}>20\,\si{\GeV}$, $|\eta|<2.3$ on both sides \\
jet--$\tau$ overlap removal & \yes\ identical & di-$\tau$ pair chosen first, then $\dR>0.4$ on both sides \\
$\ptmiss$ resolution & \yes\ tuned & Gaussian smearing to the anchor resolution ($16.5\rightarrow32.6\,\si{\GeV}$ per component) \\
lepton reconstruction efficiency & \yes\ tuned & per-$p_{T}$ scale factors (large: $1.57\rightarrow1.22$ for $e$, $1.30\rightarrow1.13$ for $\mu$) \\
di-$\tau$ mass estimator & \yes\ identical & the same covariance-free FastMTT is run on both sides \\
\bottomrule
\end{tabular}
\caption{Quantities that agree with CMS or are transferred from it.\label{tab:matched}}
\end{table}

\subsection{Residual differences, by choice or by limitation}
\label{sec:residual}

\begin{enumerate}
\item \textbf{A $\tauh$ is an AK4 jet.} This is the structural difference. In CMS a
$\tauh$ is an HPS object built from a narrow set of decay products; in Delphes it is an
$R=0.4$ jet carrying \texttt{TauTag}. Its mass, energy and direction are therefore all
wrong at source, and we correct the first two against the anchor rather than fixing the
representation. This is the subject of Section~\ref{sec:decision}.

\item \textbf{The $\tauh$ fake rate is measured on signal.} \texttt{tau\_mistag} comes
from the $b$-rich HH$\rightarrow b\bar b\tau\tau$ anchor. Fake rates are strongly
jet-flavour dependent, so this value is not appropriate for the fake-dominated
backgrounds (QCD, W+jets); those need their own measurement before they are trusted.

\item \textbf{$\ptmiss$ smearing is isotropic and flat.} Real $\ptmiss$ resolution is
anisotropic (larger along the hadronic recoil) and grows with $\Sigma E_{T}$. Ours is
round and constant, which reproduces the overall width but not its structure --- and
$\mtautau$ is now marginally \emph{broader} on Delphes than on CMS, as one would expect
from over-simplified noise.

\item \textbf{Pileup enters only through $\ptmiss$.} The card runs without pileup by
design. The jet energy \emph{resolution} is therefore untuned (only the scale is
corrected), and pileup jets cannot be produced at all. Both act on the $b\bar b$ side,
which currently agrees well; the second cannot be fixed downstream in any case.

\item \textbf{The b-jet energy scale is self-anchored.} It is corrected to its own
\texttt{GenJet} rather than to CMS. This is defensible --- CMS jet energies are
JEC-calibrated to gen --- and $m_{bb}$ agrees, but it is not the like-for-like comparison
used everywhere else.

\item \textbf{The sampled $\tauh$ mass is drawn independently} of the $\tau$ energy and
of the partner leg, whereas the true visible mass correlates with the decay mode and
hence with prong multiplicity and energy sharing.

\item \textbf{FastMTT is covariance-free} (a fixed $\ptmiss$ resolution rather than the
per-event covariance). This is run identically on both sides, so it largely cancels in
the comparison; it limits the absolute $\mtautau$ resolution, not the agreement.
\end{enumerate}

\subsection{The ten NSBI features}

\begin{table}[h]
\centering\small
\begin{tabular}{lll}
\toprule
feature & status & comment \\
\midrule
$m_{HH}$          & \yes & the $\kl$-carrying observable; agrees across $\kl=0,1,5$ \\
$m_{bb}$          & \yes & both peak near $112\,\si{\GeV}$ \\
$p_{T}^{HH}$      & \yes & \\
$\dR_{bb}$        & \yes & after symmetric overlap removal on both sides \\
$\Delta\phi_{HH}$ & \yes & \\
$p_{T}^{H1}$, $p_{T}^{H2}$ & \yes & \\
$\cos\theta^{*}$  & \yes & statistically limited but consistent \\
\midrule
$\mtautau$        & (\no) & residual $\sim$7\,\si{\GeV} high; Delphes slightly broader \\
$\dR_{\tau\tau}$  & \no  & Delphes excess at $2.2$--$3.2$ \\
\bottomrule
\end{tabular}
\caption{Agreement of the NSBI inputs with CMS at $\kl=0$ (the other $\kl$ points behave
the same). Yields: gen-matched cutflow $1.01$; analysis-level selection $1.18\times$.}
\end{table}

\paragraph{What is known about the two that disagree.} A dedicated study (Section~3,
item~6) compared the \emph{generator-level} $\dR_{\tau\tau}$ and the selection efficiency
against it, on both sides. The gen spectra agree --- so the samples are the same physics
and no generator-level explanation is available --- but the efficiencies differ: Delphes
selects $0.22$--$0.23$ of events at gen $\dR>2.5$ where CMS selects $0.15$--$0.17$
($\sim$3$\sigma$), with a weaker ($\sim$2$\sigma$) hint that Delphes also loses pairs
below $\dR<1$. Both push the normalised distribution rightward. Three explanations were
proposed and excluded by measurement: fake $\tauh$ (fake fractions are $5.7\%$ Delphes
against $5.0\%$ CMS, and the discrepancy lives entirely in the gen-matched component);
$\tau$-pair selection (making it fully symmetric moved $\dR_{\tau\tau}$ but not
$\mtautau$); and a large CMS $\tau$ energy miscalibration (the anchor response is
$0.97$--$0.99$).

\section{The decision: keep correcting the jet, or rebuild the $\tauh$?}
\label{sec:decision}

Items 1 and 6 of Section~\ref{sec:residual}, and arguably the $\dR_{\tau\tau}$
disagreement, all descend from representing a $\tauh$ as a jet. The alternative is to
rebuild it from the Delphes particle-flow candidates, which the card already writes (the
three \texttt{EFlow} collections --- tracks, photons, neutral hadrons --- kept, per the
card header, precisely for ``$\tau$ redefinition'').

\paragraph{What rebuilding would involve.} An HPS-like construction: seed from a jet,
collect charged hadrons and photons in a narrow signal cone (fixed $\dR\simeq0.1$, or the
$p_{T}$-dependent $3.0/p_{T}$ bounded at $0.05$--$0.1$), form the standard decay-mode
hypotheses (1-prong, 1-prong$+\pi^{0}$, 3-prong), take the $\tauh$ four-vector as the sum
of signal-cone constituents, and use an isolation annulus as a fake handle.

\paragraph{In favour.} Mass, energy and direction become correct \emph{by construction}
rather than by correction, and three of the current corrections --- the visible-mass map,
the CMS-anchored energy scale and the visible-$\tau$ reference --- are \emph{retired}
rather than joined by a fourth. For a phenomenological paper this is a materially
stronger statement than a jet-plus-corrections chain, and it removes a fast-simulation
limitation instead of documenting one. On cost, the implementation is contained and is
\emph{not} on the critical path: running the full $5\times10^{5}$-event samples through
the ladder and training the NSBI is the multi-day item.

\paragraph{Against, or at least uncertain.} It cannot reproduce DeepTau, so the
identification efficiency and fake-rate maps remain (they already work). It is \emph{not}
obviously the cure for $\dR_{\tau\tau}$: CMS HPS also seeds from AK4 jets, so the two
merging scales are similar and a narrow signal cone may not change which pairs are
resolved. And the current jet-based treatment now agrees with CMS on eight of ten
features and on the yield, so the marginal gain may be small.

\paragraph{The caveat we most want checked.} Delphes neutral particle-flow candidates are
\emph{calorimeter towers} ($\sim0.02\times0.02$ in the barrel), not clustered PF objects.
Since the $\pi^{0}$ content sets both the visible mass and the decay-mode assignment, an
HPS-like reconstruction on Delphes would be approximate in exactly the place that matters
most --- possibly no better than the anchor-sampled mass we already use.

\paragraph{Questions for $\tau$ expertise.}
\begin{enumerate}
\item Is tower-granular neutral reconstruction adequate for a useful decay-mode
assignment, or does it degrade the $\pi^{0}$ content enough that the rebuilt mass is no
better than the anchor-sampled one?
\item Signal cone: fixed $0.1$ or $p_{T}$-dependent, at the $p_{T}$ range that dominates
here ($20$--$60\,\si{\GeV}$)?
\item Seeding: jets (as CMS HPS does) or direct clustering of particle-flow candidates?
Would either change the efficiency for collimated pairs, $\dR_{\tau\tau}\lesssim1$?
\item Is the observed pattern --- over-efficient at gen $\dR_{\tau\tau}>2.5$, possibly
under-efficient below $1$ --- recognisable as a jet-versus-HPS effect at all, or should we
look elsewhere?
\item Is a jet-based $\tauh$ with anchor-measured corrections defensible in a
phenomenological paper, or is rebuilding the expected standard?
\item Separately: is a $1.18\times$ analysis-level yield excess acceptable for an unbinned
NSBI study, where training uses shapes and the normalisation is fitted?
\end{enumerate}

\section{Reference: the studies behind these numbers}

Each item is reproducible from the \texttt{delphes-pipeline} repository; scripts are
named where relevant.

\begin{enumerate}
\item \textbf{$\tautau$ cutflow} (\texttt{scripts/tauh\_cutflow.py}). Five stages --- gen
visible $\tauh$, reco object, ID, cleaned jets, FastMTT --- walked identically on both
samples, comparing \emph{conditional} efficiencies. Initially the FastMTT stage alone
carried the whole deficit (Delphes retained $12\%$ against CMS's $99\%$); it now agrees
at every stage, total $1.01$. Unit tests inject a loss at one known stage at a time and
assert it is attributed there and nowhere else.

\item \textbf{The FastMTT failure.} The hadronic decay prior is zero unless
$x_{\min}=(\mvis/\mtau)^{2}\le1$. A Delphes $\tau$-jet carries the AK4 jet mass
($\mathcal{O}(5\text{--}20)\,\si{\GeV}$) against $\mtau=1.777\,\si{\GeV}$, so
\emph{no} solution exists and $\mtautau$ returns NaN --- silently removing $\sim$89\% of
pairs. Fixed by assigning the anchor $\tauh$ visible mass.

\item \textbf{Why the mass must be sampled, not averaged.} The visible mass is a
one-sided \emph{lower bound} on the energy fraction, not a smearing; its contribution to
$\mvis$ is negligible. Assigning a single median therefore \emph{relaxes} a real
constraint for every leg whose true mass is higher. Worked case: a 3-prong $a_1$
($\mvis=1.26$, floor $x_{\min}=0.503$) given the median $0.77$ ($x_{\min}=0.188$) moves
from $x=0.525$ to $x=0.250$, and $\mtautau$ from $224$ to $470\,\si{\GeV}$. Over a
realistic decay-mode mixture the median inflates the width $\sim$60\% and the mean
$\sim$9\% while leaving the \emph{mode} unchanged --- so a mode-gated $Z\rightarrow\tau\tau$
candle would not catch it. The map therefore stores per-$p_{T}$ quantiles and is sampled.

\item \textbf{The reference error, twice.} A Delphes \texttt{GenJet} is also a jet and
also contains the underlying event, so it is \emph{not} the analogue of
\texttt{GenVisTau}. Profiled against a genuine visible $\tau$ (built as
$p_{\rm vis}=p_{\tau}-p_{\nu_\tau}$, valid because a hadronic decay has exactly one
neutrino), the Delphes $\tau$-jet is $17\%$ harder at $25\,\si{\GeV}$, falling to
$\sim$2\% at $250\,\si{\GeV}$. The same error appeared independently in the energy-scale
map and in the cutflow's stage-1 acceptance.

\item \textbf{Order of application.} The maps are all measured against CMS-scale
quantities, so they must be \emph{applied} at a CMS-scale $p_{T}$. Drawing the tags on
raw jets evaluated a steeply-rising efficiency $\sim$17\% too far right; applying the
energy scale first moved the analysis-level excess $1.25\rightarrow1.18$.

\item \textbf{Generator-level $\dR_{\tau\tau}$} (\texttt{scripts/tauh\_cutflow.py
--gen-dr}). Two panels: the gen $\dR$ spectrum (do the samples \emph{start} the same?)
and the selection efficiency against it (do they \emph{select} it the same?). Only the
second is a detector statement. Result quoted in Section~1.

\item \textbf{Gen-$\tau$ construction check} (\texttt{scripts/gen\_tau\_check.py}). The
CMS anchor carries both \texttt{GenPart} and \texttt{GenVisTau}, so our visible-$\tau$
construction can be scored against CMS's own definition on the same events: efficiency
$0.991$, purity $0.998$, and it reproduces CMS's stage-1 rate to $\sim$1.6\%.

\item \textbf{Generator copies.} Three of our own errors had one cause: Delphes writes
the full Pythia history, so each physical $\tau$ appears several times, while CMS NanoAOD
\texttt{GenPart} is pruned. The rule that emerged --- anything that \textbf{counts}
generator objects must collapse copies first; anything that \textbf{matches} them
geometrically need not, since copies are collinear. This is why the energy-scale maps,
the re-tagging and the overlay were never affected, while a stage-1 \emph{count} was.
\end{enumerate}

\paragraph{Reproducing.} Derive the maps, then run the ladder and the comparison:
\begin{verbatim}
python -m delphes_pipeline.tuning.derive_maps --config config.v1.yml \
    --output cards/tuning/maps_v1.json
python -m delphes_pipeline.validation.run_validation --config config.v1.yml
python scripts/tauh_cutflow.py --config config.v1.yml [...] --gen-dr plots/gen_dr
python scripts/nsbi_overlay.py --config config.v1.yml [...] \
    --tautau-only --diagnostics
\end{verbatim}

\end{document}
